%% LyX 2.4.4 created this file.  For more info, see https://www.lyx.org/.
%% Do not edit unless you really know what you are doing.
\documentclass[english]{article}
\usepackage[T1]{fontenc}
\usepackage[latin9]{inputenc}
\usepackage{units}
\usepackage{enumitem}
\usepackage{amsmath}
\usepackage{amsthm}
\usepackage{graphicx}
\usepackage{geometry}
\geometry{verbose,tmargin=1in,bmargin=1in,lmargin=1in,rmargin=1in}

\makeatletter
%%%%%%%%%%%%%%%%%%%%%%%%%%%%%% Textclass specific LaTeX commands.
\numberwithin{equation}{section}
\numberwithin{figure}{section}
\newlength{\lyxlabelwidth}      % auxiliary length 

%%%%%%%%%%%%%%%%%%%%%%%%%%%%%% User specified LaTeX commands.
\usepackage{fancyhdr}
\usepackage{lastpage}
\pagestyle{fancy}
\usepackage{enumitem}
\usepackage{ifthen}
\everymath=\expandafter{\the\everymath\displaystyle}
%\DeclareMathSizes{10}{10}{10}{10}
\usepackage{multicol}
% Added by lyx2lyx
\setlength{\parskip}{\smallskipamount}
\setlength{\parindent}{0pt}

\makeatother

\usepackage{babel}
\begin{document}
\renewcommand{\author}{Dr.\,James Ashe}

\newcommand{\classNum}{232}

\newcommand{\classSec}{C}

\newcommand{\examType}{EXAM}

\renewcommand{\date}{October 15, 2012}

%%First page's header%%%%%%%%%%%%%%%%%%%%%%
\fancypagestyle{firststyle} %%header settings for first pagr
{
	\fancyhead[L]{MTH \classNum{}(\classSec{}) \author{}\\
	\textbf{\examType{}} \date{}}
	\fancyhead[C]{\textbf{Name:}}
	\setlength{\headsep}{14pt}
}
%%%%%%%%%%%%%%%%%%%%%%%%%%%%%%%%%%%%%%%%%%

%%Next page's header%%%%%%%%%%%%%%%%%%%%%%%
\setlist{nolistsep}
\lhead{\classNum{}(\classSec{})} 
\chead{\textbf{Name:}} 
\cfoot{\thepage{}~of~\pageref{LastPage}} 
\setlength{\headsep}{5pt} 
%%%%%%%%%%%%%%%%%%%%%%%%%%%%%%%%%%%%%%%%%%%

%%Various settings; see the preamble for other settings%%%%%
%%\setenumerate[0]{leftmargin=12pt,itemindent=0pt,labelwidth=0pt}
\renewcommand{\labelenumi}{\textbf{\arabic{enumi}.}}
\renewcommand{\labelenumii}{\textbf{\alph{enumii}.}}
\thispagestyle{firststyle}
%%%%%%%%%%%%%%%%%%%%%%%%%%%%%%%%%%%%%%%%%%

\textbf{Justify your answers and show your work.}

\emph{Complete the following steps for $f(x)=24-2x$, interval $[0,6]$,
and $n=3$.}
\begin{enumerate}[resume]
\item Sketch the graph of the function on the given interval and illustrate
the \emph{right Riemann sums }using $n$ subintervals.\vspace{2.5in}
\item Calculate the \emph{right Riemann sums }for the given interval and
number of subintervals, $n$.\vfill
\end{enumerate}
\emph{The figure shows the area of regions bounded by the graph of
$f(x)$. Evaluate the following integrals.}

\begin{multicols}{2}
\begin{enumerate}[resume]
\item ${\displaystyle \int_{0}^{8}f}(x)\,dx$\vfill
\item ${\displaystyle \int_{0}^{5}f(x)\,dx}-{\displaystyle \int_{8}^{0}f(x)\,dx}$\vfill\columnbreak
\end{enumerate}
\includegraphics[width=3in]{D:/home/jim/JamesAshe/JCSU/Calc_2_MTH_232/images/Exam_1_area_under_graph.svg}

\end{multicols}\newpage

\emph{Evaluate the following indefinite integrals and definite integrals.}
\begin{enumerate}[resume]
\item [\textbf{5.}]\setcounter{enumi}{5}$f(x)=9x^{8}$

\begin{enumerate}
\item ${\displaystyle \int9x^{8}\,dx}$\vfill
\item ${\displaystyle \int_{0}^{1}9x^{8}\,dx}$\vfill
\end{enumerate}
\item $f(x)={\displaystyle x^{2}-8x+1}$

\begin{enumerate}
\item $\int(x^{2}-8x+1)\,dx$\vfill
\end{enumerate}
\item $f(x)=\cos3x$

\begin{enumerate}
\item ${\displaystyle \int\cos3x\,dx}$\vfill
\item ${\displaystyle \int_{-\nicefrac{\pi}{3}}^{\nicefrac{\pi}{3}}\cos3x\,dx}$\vfill\newpage
\end{enumerate}
\item $f(x)=\frac{x-\sqrt{x}}{x^{3}}$

\begin{enumerate}
\item ${\displaystyle \int\frac{x-\sqrt{x}}{x^{3}}}\,dx$\vfill
\end{enumerate}
\item $f(x)=x^{2}(3x^{3}+1)^{4}$

\begin{enumerate}
\item ${\displaystyle \int x^{2}(3x^{3}+1)^{4}\,dx}$\vfill
\item ${\displaystyle \int_{1}^{2}x^{2}(3x^{3}+1)^{4}\,dx}$\vfill
\end{enumerate}
\item $v(\theta)={\displaystyle \cos\theta\sin^{2}\theta}$

\begin{enumerate}
\item ${\displaystyle \int\cos\theta\sin^{2}\theta\,d\theta}$\vfill\newpage
\end{enumerate}
\item ${\displaystyle g(t)=\frac{t}{\sqrt{4-t^{2}}}}$

\begin{enumerate}
\item ${\displaystyle \int\frac{t}{\sqrt{4-t^{2}}}\,dt}$\vspace{1.5in}
\item ${\displaystyle \int_{2}^{4}\frac{t}{\sqrt{4-t^{2}}}\,dt}$\vspace{1.25in}
\end{enumerate}
\item A hiking trail has an elevation given by $h(x)=40x^{3}-300x^{2}+200x+4000$
where $h$ measures feet above sea level and $x$ represents miles
along the trail with $0\leq x\leq5$. What is the average elevation
along the trail?\vfill
\item A bowling ball rolls down a ramp at a speed of $v(t)=-3t^{2}$ feet
per second (where $t$ is seconds). The distance of the ball from
the top of the ramp is $4$ feet after $2$ seconds. Find the function
giving the distance of the ball from the top of the ramp at $t$ seconds.\vfill
\end{enumerate}

\end{document}
